\chapter{CA-125 (Cancer Antigen 125)}

\section*{What Is This Test?}
The CA-125 test measures the concentration of cancer antigen 125 (CA-125, MUC16) in the blood. CA-125 is a high-molecular-weight (approximately 200--300 kDa) glycoprotein of the mucin family (MUC16). It is expressed by epithelial cells of the female reproductive tract, the peritoneum, pleura, and pericardium, and by mesothelial cells. It is most strongly associated with epithelial ovarian cancer (EOC), especially the serous subtype, and is also elevated in other gynecologic cancers (endometrial, fallopian tube, primary peritoneal), breast, pancreatic, gastric, colorectal, and lung cancers. CA-125 is present in high concentrations in amniotic fluid and is elevated in many benign conditions (endometriosis, pelvic inflammatory disease, pregnancy, uterine fibroids, ovarian cysts, ascites, pleuritis, hepatitis, pancreatitis). It is used primarily for (1) monitoring response to treatment and detecting recurrence in patients with a known diagnosis of epithelial ovarian cancer; and (2) evaluating an adnexal mass, although its role in screening asymptomatic women is not established. CA-125 is NOT a screening test for ovarian cancer in the general population because of low sensitivity and specificity. The Risk of Ovarian Malignancy Algorithm (ROMA) combines CA-125 with HE4 (human epididymis protein 4) and menopausal status to stratify the risk of malignancy in women with a pelvic mass.

\section*{Alternative Names}
\begin{itemize}
  \item Cancer Antigen 125
  \item CA-125
  \item CA125
  \item MUC16 (Cancer Antigen)
  \item Ovarian Cancer Antigen
  \item Tumor Marker CA-125
\end{itemize}
\section*{Principle}
CA-125 is measured by immunoassay, typically electrochemiluminescent immunoassay (ECLIA) or chemiluminescent microparticle immunoassay (CMIA) on automated analyzers. The antibody used (OC125) is a monoclonal antibody raised against the OVCA 433 ovarian cancer cell line. CA-125 is quantified in units per milliliter (U/mL or kU/L). The assay has high reproducibility within a laboratory, but values from different laboratories can vary by 10--20\% because different assays use different antibodies and calibrators. Serial CA-125 measurements from the same laboratory are essential for meaningful trending. CA-125 is elevated in serum when tumor burden is large or when mesothelial/peritoneal surfaces are irritated. It is not a direct measure of tumor volume and can be normal in early-stage ovarian cancer (up to 50\% of stage I tumors have normal CA-125). It is also elevated in many nonmalignant conditions that involve mesothelial or peritoneal surfaces (endometriosis, ascites, pleuritis, peritonitis, hepatic cirrhosis with ascites).

\section*{History}
CA-125 was discovered in 1981 by Robert Bast Jr. and colleagues at the Dana-Farber Cancer Institute. They developed a monoclonal antibody (OC125) that reacted with a surface antigen of the OVCA 433 serous ovarian cancer cell line. The antigen was named CA-125, and the first radioimmunoassay was published in 1983 \cite{bast1981reactivity}\cite{bast1983radioimmunoassay}. CA-125 became the standard tumor marker for epithelial ovarian cancer. In the 1990s, the UKCTOCS trial (UK Collaborative Trial of Ovarian Cancer Screening) and other studies evaluated CA-125-based screening. The results did not demonstrate a mortality benefit from screening the general population, and CA-125 is not recommended for screening. The Risk of Malignancy Index (RMI) and later the Risk of Ovarian Malignancy Algorithm (ROMA) were developed to improve the evaluation of adnexal masses. The UKCTOCS trial (2016) evaluated multimodal screening with CA-125 and transvaginal ultrasound and found no significant mortality reduction \cite{jacobs2016ovarian}. CA-125 remains the standard of care for monitoring treatment response and recurrence surveillance in ovarian cancer, and HE4 added to CA-125 improves specificity.

\section*{Why This Test Is Ordered}
\begin{itemize}
  \item Monitoring response to treatment in patients with a known diagnosis of epithelial ovarian cancer: CA-125 is measured before and during chemotherapy. A fall of >50\% in CA-125 (e.g., from 500 to <250 U/mL) correlates with a response to treatment (Rustin criteria). A rising CA-125 during treatment suggests disease progression
  \item Surveillance for recurrence after primary treatment: in patients with elevated CA-125 at diagnosis, a rising CA-125 is often the first sign of recurrence, preceding radiologic findings by 2--6 months. The OV05/EORTC 55955 trial (2010) showed that immediate treatment based on rising CA-125 alone (without symptoms or imaging evidence) did not improve survival, so the decision to treat a rising CA-125 is individualized
  \item Evaluation of a pelvic/adnexal mass: CA-125 is measured in women with a pelvic mass. In premenopausal women, CA-125 >200 U/mL is highly suspicious for malignancy; in postmenopausal women, >65--100 U/mL is suspicious. CA-125 is combined with HE4 in the ROMA score. In premenopausal women, benign conditions (endometriosis, fibroids, PID) frequently elevate CA-125, so elevation alone is not diagnostic
  \item Assessment of ascites or pleural effusion of unknown origin: CA-125 is often markedly elevated (>1000 U/mL) in the presence of large-volume ascites or pleural effusions, regardless of whether the cause is benign (cirrhosis, CHF) or malignant. A very high CA-125 with benign causes is a well-recognized pitfall
  \item Research and clinical trials: CA-125 kinetics are used as a surrogate endpoint (Rustin criteria) in some trials of ovarian cancer therapy
\end{itemize}

\section*{Primary vs Secondary Test}
CA-125 is a primary test for monitoring ovarian cancer treatment response and recurrence surveillance. It is a secondary/supplementary test for the evaluation of a pelvic mass (combined with HE4 and imaging) and is NOT recommended for screening the general population.

\section*{How This Test Is Performed}
A healthcare professional draws blood from a vein into a serum separator tube (gold-top). The sample is centrifuged and analyzed by immunoassay on automated analyzers. Results are typically available within 1--2 hours. Serial values for monitoring are most reliable when measured with the same assay at the same laboratory.

\section*{What Preparation Is Needed}
No fasting is required. Menstrual cycle phase affects CA-125 in premenopausal women (peak during menstruation). If possible, schedule the test 5--7 days after the start of the last menstrual period for a baseline assessment. Document all current medications, especially antibiotics, hormonal therapy, and chemotherapy.

\section*{Contraindications and Precautions}
\begin{itemize}
  \item CA-125 is NOT recommended for ovarian cancer screening in the general population. Low sensitivity (many early-stage cancers have normal CA-125) and low specificity (many benign conditions cause elevation) result in unacceptable false-positive and false-negative rates. Screening with CA-125 + transvaginal ultrasound did not reduce ovarian cancer mortality (UKCTOCS)
  \item Many benign conditions elevate CA-125: menstruation, pregnancy (1st trimester), endometriosis, adenomyosis, uterine fibroids, ovarian cysts, pelvic inflammatory disease (PID), tubo-ovarian abscess, peritonitis, pleuritis, pericarditis, ascites (any cause), cirrhosis, hepatitis, pancreatitis, tuberculosis, recent abdominal surgery, and postoperative states
  \item Ascites and large pleural effusions cause CA-125 elevation independent of malignancy: the mesothelial irritation itself raises CA-125. A patient with cirrhosis and ascites may have CA-125 >500--1000 U/mL without any malignancy
  \item CA-125 is falsely normal in up to 50\% of stage I ovarian cancers and in some mucinous, clear cell, and low-grade serous subtypes. A normal CA-125 does NOT exclude ovarian cancer
  \item CA-125 rises during chemotherapy cycles (transient, due to mesothelial irritation by cytotoxic drugs); a single elevated value during treatment does not necessarily indicate progression. Serial measurement (same assay, same lab) with clinical and radiologic correlation is required
  \item Premenopausal women normally have higher baseline CA-125 than postmenopausal women. Reference ranges differ by menopausal status and by assay
  \item Patients with the BRCA1/BRCA2 mutations and hereditary breast and ovarian cancer syndrome: CA-125 + transvaginal ultrasound every 6 months is an option for surveillance, but it has not been proven to reduce mortality. Prophylactic bilateral salpingo-oophorectomy remains the definitive risk-reducing strategy
\end{itemize}
\section*{Risks}
Standard venipuncture risks: mild pain, bruising, or hematoma at the puncture site.

\section*{What Happens After the Test}
Results are available within 1--2 hours. A single CA-125 value must be interpreted in clinical context. For monitoring, the trend matters more than any single value. For a pelvic mass, CA-125 is combined with HE4 (ROMA), transvaginal ultrasound, and menopausal status. If CA-125 is elevated and malignancy is confirmed, serial measurements guide treatment response and detect recurrence.

\section*{Understanding Results}

\subsection*{Below Normal (Normal Range)}
\begin{itemize}
  \item \textbf{CA-125 <35 U/mL (premenopausal) or <20 U/mL (postmenopausal) with a pelvic mass:} Low likelihood of epithelial ovarian cancer, but does not exclude it (mucinous and clear cell tumors may have normal CA-125). Combined with a benign-appearing ultrasound, this supports conservative management or surveillance
  \item \textbf{Normal CA-125 during ovarian cancer follow-up:} No evidence of recurrence. However, recurrence can occur without a rise in CA-125, particularly in mucinous and clear cell subtypes. Imaging and symptoms guide further evaluation
  \item \textbf{Normal CA-125 with ascites:} Unlikely to be epithelial ovarian cancer; consider peritoneal mesothelioma, lymphoma, or benign causes
\end{itemize}

\subsection*{Above Normal}
\begin{itemize}
  \item \textbf{CA-125 >200 U/mL in a premenopausal woman with a pelvic mass:} Strongly suspicious for epithelial ovarian cancer (though endometriosis, PID, and fibroids can produce this range). Prompt gynecologic oncology referral is warranted
  \item \textbf{CA-125 >65 U/mL in a postmenopausal woman with a pelvic mass:} Suspicious for ovarian cancer. Postmenopausal women rarely have benign CA-125 elevations, so the threshold is lower. MRI, CT, or transvaginal ultrasound, and gynecologic oncology consultation are recommended
  \item \textbf{CA-125 >1000 U/mL:} Essentially diagnostic of advanced ovarian cancer (or rarely, advanced non-ovarian cancers with peritoneal carcinomatosis, or large-volume ascites from benign causes). In a symptomatic woman with ascites and an adnexal mass, this level strongly supports the diagnosis of advanced epithelial ovarian cancer
  \item \textbf{Rising CA-125 during follow-up (e.g., doubling from 20 to 40 U/mL):} Raises suspicion of recurrence. CA-125 >2x the upper limit of normal (ULN) or >35 U/mL on two consecutive samples 2--4 weeks apart, in a patient previously treated for ovarian cancer, is considered evidence of biochemical recurrence. Whether to treat on the basis of CA-125 alone is an individualized decision (OV05/EORTC 55955 showed no survival benefit from immediate treatment)
\end{itemize}

\section*{Normal Results}
\begin{itemize}
  \item \textbf{Premenopausal women:} 0--35 U/mL (typical reference range; values up to 65 U/mL may occur in menstruation, ovulation, and benign gynecologic disease)
  \item \textbf{Postmenopausal women:} 0--20 U/mL (typical reference range)
  \item \textbf{Men:} 0--35 U/mL (CA-125 is rarely measured in men; normal values are low)
\end{itemize}

\section*{Follow-up Tests}
\begin{itemize}
  \item \textbf{HE4 (human epididymis protein 4):} More specific than CA-125 for epithelial ovarian cancer and less affected by benign gynecologic disease. HE4 is combined with CA-125 in the Risk of Ovarian Malignancy Algorithm (ROMA) score to stratify the risk of malignancy in women with a pelvic mass. HE4 is elevated in a smaller proportion of benign conditions than CA-125
  \item \textbf{Transvaginal ultrasound (TVUS):} First-line imaging for a pelvic mass. Morphology (complex cyst, solid component, mural nodule, septations, ascites) combined with CA-125 defines the Risk of Malignancy Index (RMI). CT abdomen/pelvis is used for staging and surgical planning
  \item \textbf{CA-125 monitoring (serial):} Every 2--3 months during first-line chemotherapy; every 3--6 months during surveillance for the first 2--3 years, then every 6--12 months. Trends, not absolute values, guide management
  \item \textbf{CT, MRI, or PET/CT:} For staging, assessment of disease burden, and detection of recurrence (especially retroperitoneal lymphadenopathy, peritoneal carcinomatosis, and liver/pleural involvement)
  \item \textbf{BRCA1/BRCA2 genetic testing:} In women with epithelial ovarian cancer (especially high-grade serous), testing for germline and somatic BRCA mutations guides PARP inhibitor therapy (olaparib) and informs risk-reduction planning for family members
  \item \textbf{CA-19-9 and CEA:} May be elevated in mucinous and clear cell ovarian cancer subtypes; used in combination with CA-125 when these histologies are suspected. Endometrial, fallopian tube, and primary peritoneal cancers are also monitored with CA-125
\end{itemize}

\section*{References}
\bibliographystyle{unsrtnat}
\bibliography{references}

\clearpage
\section*{Regulatory Status and Authorization Date}
This chapter describes a test category, not a specific commercial product. FDA, CDSCO, EU IVDR, and MHRA approval or registration dates therefore cannot be assigned without the assay manufacturer, product name, instrument, and intended use. For laboratory-developed tests, an individual agency approval date may not exist. See the Preface for the interpretation of regulatory status.

\clearpage
